\subsection{Kernels and Nonlinear Methods}
% RKHS intuition, kernel ridge regression, SVMs, Gaussian processes as learning models.

\begin{remark} \hlt{Nonlinear Models}\\
To represent nonlinear functions of $\bm{x}$, apply linear model to transformed input $\phi(\bm{x})$ where $\phi$ is a nonlinear transformation. Kernel trick may be applied to obtain nonlinear learning algorithm. To choose mapping $\phi$:
\begin{enumerate}[label=\roman*.]
\setlength{\itemsep}{0pt}
\item Choose generic $\phi$, such as that used by kernel machines based on RBF kernel.\\
If $\phi(\bm{x})$ is of high enough dimension, can find enough capacity to fit training set, but generalisations to test set remains poor. Mappings are based on principle of local smoothness and do not encode enough prior information to solve advanced problems.
\item Manually engineer $\phi$. Requires decades of human effort for each task, and practitioners specialising in different domains (i.e, speech recognition, computer vision) with little transfer between domains.
\item Learn $\phi$ through deep learning. Model is $y = f(\bm{x}; \bm{\theta}, \bm{w}) = \phi(\bm{x}; \bm{\theta})^T \bm{w}$, where parameters $\bm{\theta}$ can be used to learn $\phi$ from broad class of functions, and parameters $\bm{w}$ that map from $\phi(\bm{x})$ to desired output.\\
Do not require training problem to be convex, and only require finding the right general function.
\end{enumerate}
\end{remark}
